\documentclass[12pt,letterpaper]{article}
\usepackage{graphicx,textcomp}
\usepackage{natbib}
\usepackage{setspace}
\usepackage{fullpage}
\usepackage{color}
\usepackage[reqno]{amsmath}
\usepackage{amsthm}
\usepackage{fancyvrb}
\usepackage{amssymb,enumerate}
\usepackage[all]{xy}
\usepackage{endnotes}
\usepackage{lscape}
\newtheorem{com}{Comment}
\usepackage{float}
\usepackage{hyperref}
\newtheorem{lem} {Lemma}
\newtheorem{prop}{Proposition}
\newtheorem{thm}{Theorem}
\newtheorem{defn}{Definition}
\newtheorem{cor}{Corollary}
\newtheorem{obs}{Observation}
\usepackage[compact]{titlesec}
\usepackage{dcolumn}
\usepackage{tikz}
\usetikzlibrary{arrows}
\usepackage{multirow}
\usepackage{xcolor}
\newcolumntype{.}{D{.}{.}{-1}}
\newcolumntype{d}[1]{D{.}{.}{#1}}
\definecolor{light-gray}{gray}{0.65}
\usepackage{url}
\usepackage{listings}
\usepackage{color}

\definecolor{codegreen}{rgb}{0,0.6,0}
\definecolor{codegray}{rgb}{0.5,0.5,0.5}
\definecolor{codepurple}{rgb}{0.58,0,0.82}
\definecolor{backcolour}{rgb}{0.95,0.95,0.92}

\lstdefinestyle{mystyle}{
	backgroundcolor=\color{backcolour},   
	commentstyle=\color{codegreen},
	keywordstyle=\color{magenta},
	numberstyle=\tiny\color{codegray},
	stringstyle=\color{codepurple},
	basicstyle=\footnotesize,
	breakatwhitespace=false,         
	breaklines=true,                 
	captionpos=b,                    
	keepspaces=true,                 
	numbers=left,                    
	numbersep=5pt,                  
	showspaces=false,                
	showstringspaces=false,
	showtabs=false,                  
	tabsize=2
}
\lstset{style=mystyle}
\newcommand{\Sref}[1]{Section~\ref{#1}}
\newtheorem{hyp}{Hypothesis}


\title{Problem Set 4}
\date{Due: November 27, 2025}
\author{Applied Stats/Quant Methods 1}


\begin{document}
	\maketitle
	\section*{Instructions}
	\begin{itemize}
		\item Please show your work! You may lose points by simply writing in the answer. If the problem requires you to execute commands in \texttt{R}, please include the code you used to get your answers. Please also include the \texttt{.R} file that contains your code. If you are not sure if work needs to be shown for a particular problem, please ask.
		\item Your homework should be submitted electronically on GitHub.
		\item This problem set is due before 23:59 on Thursday November 27, 2025. No late assignments will be accepted.
	\end{itemize}



	\vspace{.5cm}
\section*{Question 1: Economics}
\vspace{.25cm}
\noindent 	
In this question, use the \texttt{prestige} dataset in the \texttt{car} library. First, run the following commands:

\begin{verbatim}
install.packages(car)
library(car)
data(Prestige)
help(Prestige)
\end{verbatim} 


\noindent We would like to study whether individuals with higher levels of income have more prestigious jobs. Moreover, we would like to study whether professionals have more prestigious jobs than blue and white collar workers.

\newpage
\begin{enumerate}
	
	\item [(a)]
	Create a new variable \texttt{professional} by recoding the variable \texttt{type} so that professionals are coded as $1$, and blue and white collar workers are coded as $0$ (Hint: \texttt{ifelse}).
	
	\textbf{Answer:}
	We want \texttt{professional} = 1 if type == 'prof', and 0 otherwise.
	
	\vspace{1em}
	Computed in \texttt{R} as:
	\lstinputlisting[language=R, firstline=46, lastline=47]{PS04_JL.R} 
	
	
	\item [(b)]
	Run a linear model with \texttt{prestige} as an outcome and \texttt{income}, \texttt{professional}, and the interaction of the two as predictors (Note: this is a continuous $\times$ dummy interaction.)
	
	\textbf{Answer:}
	We regress \texttt{prestige} on \texttt{income}, \texttt{professional}, and their interaction
	
	\begin{table}[!htbp] \centering 
		\caption{Regression of Prestige on Income, Professional, and Interaction} 
		\label{} 
		\begin{tabular}{@{\extracolsep{5pt}}lc} 
			\\[-1.8ex]\hline 
			\hline \\[-1.8ex] 
			& \multicolumn{1}{c}{\textit{Dependent variable:}} \\ 
			\cline{2-2} 
			\\[-1.8ex] & prestige \\ 
			\hline \\[-1.8ex] 
			income & 0.003$^{***}$ \\ 
			& (0.0005) \\ 
			& \\ 
			professional & 37.781$^{***}$ \\ 
			& (4.248) \\ 
			& \\ 
			income:professional & $-$0.002$^{***}$ \\ 
			& (0.001) \\ 
			& \\ 
			Constant & 21.142$^{***}$ \\ 
			& (2.804) \\ 
			& \\ 
			\hline \\[-1.8ex] 
			Observations & 98 \\ 
			R$^{2}$ & 0.787 \\ 
			Adjusted R$^{2}$ & 0.780 \\ 
			Residual Std. Error & 8.012 (df = 94) \\ 
			F Statistic & 115.878$^{***}$ (df = 3; 94) \\ 
			\hline 
			\hline \\[-1.8ex] 
			\textit{Note:}  & \multicolumn{1}{r}{$^{*}$p$<$0.1; $^{**}$p$<$0.05; $^{***}$p$<$0.01} \\ 
		\end{tabular} 
	\end{table} 
	\vspace{1em}
	Computed in \texttt{R} as:
	\lstinputlisting[language=R, firstline=49, lastline=51]{PS04_JL.R} 
	\newpage

	\item [(c)]
	Write the prediction equation based on the result.
	
	\textbf{Answer:}
	\[
	\widehat{\texttt{prestige}} = 21.142 + 0.003 \cdot \texttt{income} + 37.781 \cdot \texttt{professional} - 0.002 \cdot (\texttt{income} \times \texttt{professional})
	\]
	

	\item [(d)]
	Interpret the coefficient for \texttt{income}.
	
	\textbf{Answer:}  
	The coefficient for \texttt{income} represents the expected change in prestige score for non-professional occupations (\texttt{professional} == 0) when income increases by one unit. In this case, a one-dollar increase in income raises prestige by 0.003 points. For a more meaningful scale, a \$1{,}000 increase in income is associated with a 3-point increase in prestige among non-professional jobs.
	
	\item [(e)]
	Interpret the coefficient for \texttt{professional}.
	
	\textbf{Answer:}  
	The coefficient for \texttt{professional} represents the difference in prestige between professional and non-professional occupations when income is zero. Although \texttt{income} == 0 is not realistic, the interpretation is that professionals start with a substantially higher baseline prestige score compared to non-professionals, before accounting for income.

	\item [(f)]
	What is the effect of a \$1,000 increase in income on prestige score for professional occupations? In other words, we are interested in the marginal effect of income when the variable \texttt{professional} takes the value of $1$. Calculate the change in $\hat{y}$ associated with a \$1,000 increase in income based on your answer for (c).
	
	\textbf{Answer:} 
	For professional occupations, the marginal effect ($\text{ME}$) of income is obtained by taking the partial derivative of the prediction equation with respect to \texttt{income}. This derivative simplifies to the sum of the \texttt{income} coefficient and the interaction coefficient. In numerical terms, $0.003 + (-0.002) = 0.001$. Therefore, a \$1{,}000 increase in income raises the prestige score by $0.001 \times 1000 = 1$ point for professionals. Put differently, while income generally has a positive effect on \texttt{prestige}, this effect is reduced by 0.002 for individuals in professional occupations.
	\[
	\text{ME}_{\texttt{income}} = \frac{\partial \widehat{\texttt{prestige}}}{\partial \texttt{income}} = 0.003 - 0.002 \cdot \texttt{professional}
	\]
	\newpage
	
	\item [(g)]
	What is the effect of changing one's occupations from non-professional to professional when her income is \$6,000? We are interested in the marginal effect of professional jobs when the variable \texttt{income} takes the value of $6,000$. Calculate the change in $\hat{y}$ based on your answer for (c).
	
	\textbf{Answer:}  
	The marginal effect ($\text{ME}$) of being in a professional occupation is obtained by taking the partial derivative of the prediction equation with respect to \texttt{professional}. This derivative simplifies to the sum of the \texttt{professional} coefficient and the interaction coefficient multiplied by \texttt{income}. In numerical terms, $37.781 + (-0.002 \times 6000) = 25.781$. Therefore, at an income of \$6{,}000, switching from a non-professional to a professional occupation increases the prestige score by 25 points. In other words, the positive effect of being a professional is diminished as income rises, because the interaction term reduces the gap between professionals and non-professionals at higher income levels.	
	\[
	\text{ME}_{\texttt{professional}} = \frac{\partial \widehat{\texttt{prestige}}}{\partial \texttt{professional}} = 37.781 + (-0.002) \cdot \texttt{income}
	\]
	
\end{enumerate}

\newpage

\section*{Question 2: Political Science}
\vspace{.25cm}
\noindent 	Researchers are interested in learning the effect of all of those yard signs on voting preferences.\footnote{Donald P. Green, Jonathan	S. Krasno, Alexander Coppock, Benjamin D. Farrer,	Brandon Lenoir, Joshua N. Zingher. 2016. ``The effects of lawn signs on vote outcomes: Results from four randomized field experiments.'' Electoral Studies 41: 143-150. } Working with a campaign in Fairfax County, Virginia, 131 precincts were randomly divided into a treatment and control group. In 30 precincts, signs were posted around the precinct that read, ``For Sale: Terry McAuliffe. Don't Sellout Virgina on November 5.'' \\

Below is the result of a regression with two variables and a constant.  The dependent variable is the proportion of the vote that went to McAuliff's opponent Ken Cuccinelli. The first variable indicates whether a precinct was randomly assigned to have the sign against McAuliffe posted. The second variable indicates
a precinct that was adjacent to a precinct in the treatment group (since people in those precincts might be exposed to the signs).  \\

\vspace{.5cm}
\begin{table}[!htbp]
	\centering 
	\textbf{Impact of lawn signs on vote share}\\
	\begin{tabular}{@{\extracolsep{5pt}}lccc} 
		\\[-1.8ex] 
		\hline \\[-1.8ex]
		Precinct assigned lawn signs  (n=30)  & 0.042\\
		& (0.016) \\
		Precinct adjacent to lawn signs (n=76) & 0.042 \\
		&  (0.013) \\
		Constant  & 0.302\\
		& (0.011)
		\\
		\hline \\
	\end{tabular}\\
	\footnotesize{\textit{Notes:} $R^2$=0.094, N=131}
\end{table}

\vspace{.5cm}
\begin{enumerate}
	\item [(a)] Use the results from a linear regression to determine whether having these yard signs in a precinct affects vote share (e.g., conduct a hypothesis test with $\alpha = .05$).
	
	\textbf{Answer:}  
	We test whether the coefficient for \texttt{assigned} (precincts with yard signs) differs significantly from zero. The estimated coefficient is $0.042$ with a standard error of $0.016$. The test statistic is:
	\[
	t = \frac{0.042}{0.016} = 2.625
	\]
	Since $|t| > 1.96$, we reject the null hypothesis at the $\alpha = 0.05$ level. We conclude that having yard signs in a precinct has a statistically significant positive effect on Cuccinelli’s vote share. The 95\% confidence interval is:
	\[
	0.042 \pm 1.96 \times 0.016 \Rightarrow [0.011,\; 0.073]
	\]
	\newpage
	
	\item [(b)]  Use the results to determine whether being next to precincts with these yard signs affects vote share (e.g., conduct a hypothesis test with $\alpha = .05$).
	
	\textbf{Answer:}  
	We test whether the coefficient for \texttt{adjacent} (precincts next to those with yard signs) differs significantly from zero. The estimated coefficient is $0.042$ with a standard error of $0.013$. The test statistic is:
	\[
	t = \frac{0.042}{0.013} \approx 3.231
	\]
	Since $|t| > 1.96$, we reject the null hypothesis at the $\alpha = 0.05$ level. We conclude that being adjacent to a precinct with yard signs also has a statistically significant positive effect on Cuccinelli’s vote share. The 95\% confidence interval is:
	\[
	0.042 \pm 1.96 \times 0.013 \Rightarrow [0.017,\; 0.067]
	\]
	
	\item [(c)] Interpret the coefficient for the constant term substantively.
	
	\textbf{Answer:}  
	The constant term is estimated at $0.302$ with a standard error of $0.011$. This represents the expected vote share for Cuccinelli in precincts that had no yard signs and were not adjacent to any precincts with signs. Substantively, it suggests that in the absence of any sign exposure, Cuccinelli’s baseline vote share was approximately 30.2\%.
	
	\item [(d)] Evaluate the model fit for this regression.  What does this	tell us about the importance of yard signs versus other factors that are not modeled?
	
	\textbf{Answer:}  
	The model has an $R^2$ of $0.094$, indicating that the two predictors (\texttt{assigned} and \texttt{adjacent}) explain 9.4\% of the variation in Cuccinelli’s vote share across precincts. While the effects of yard signs and adjacency are statistically significant and nontrivial in magnitude (about 4.2 percentage points each), the relatively low $R^2$ suggests that most of the variation in vote share is driven by other factors not included in the model such as demographics, partisanship, turnout, or campaign outreach, conceivably. Yard signs matter, but they account for only a small portion of the overall variation in voting behaviour.
	
\end{enumerate}  


\end{document}
